\section*{章末問題}
\begin{enumerate}[label=\textbf{\thechapter.\arabic*},leftmargin=*]
\item\label{問:奇偶の混合}
次の定積分を求めよ。
\[
\int_{-2}^2\left(\frac{x^3}{1+x^2}+2x^2-\cos x\right)\dd x.
\]
\item\label{問:一次重み反転}
反転を利用して，$\int_0^\pi x\sin^2x\dd x$ を求めよ。
\item\label{問:三角相補}
$\int_0^{\frac{\pi}{2}}\frac{\cos^6x}{\sin^6x+\cos^6x}\dd x$ を求めよ。分母が零にならないことも確かめよ。
\item\label{問:周期絶対値}
$\int_{\frac{\pi}{3}}^{\frac{10\pi}{3}}|\cos x|\dd x$ と
$\int_{\frac{\pi}{3}}^{\frac{10\pi}{3}}\cos x\dd x$ をそれぞれ求め，値が異なる理由を述べよ。
\item\label{問:逆数四乗}
$\int_{\frac14}^{4}\frac{1}{x(1+x^4)}\dd x$ を，逆数の置換を用いて求めよ。
\item\label{問:指数相補}
$\int_{-1}^1\frac{3x^2+1}{1+\e^{x^5}}\dd x$ を求めよ。
\item\label{問:中心を移す相補}
$\int_2^6\frac{(x-4)^2+2}{1+\e^{x-4}}\dd x$ を求めよ。
\item\label{問:二次重みの限界}
$I=\int_0^\pi x^2\sin x\dd x$ とする。反転によって
$I=\int_0^\pi(\pi-x)^2\sin x\dd x$ と書けるが，この二式を加えるだけで $I$ は決まらないことを説明せよ。その後，部分積分を用いて $I$ を求めよ。
\end{enumerate}

\smallskip\noindent 解答は\pageref{解答:第7章}ページ。
